\chapter{Біль: сигнал тривоги}
\chaptermark{Біль}
\label{sec:pain}

\section{Тривога, а не вимірювання}

Усі датчики розділу~\ref{sec:sensors} вимірювали світ: скільки світла, якої висоти звук, який тиск. Біль улаштований інакше. Він повідомляє не «що там», а «небезпечно, кидай усе». Для інженера це не вимірювальний канал, а \emph{сигнал тривоги}: лінія з найвищим пріоритетом, яка має пробитися крізь будь-яку поточну роботу машини.

Від вимірювального каналу сигналізація відрізняється трьома властивостями, і всі три біль має. По-перше, його важко приглушити звиканням: датчики болю звикають набагато слабше за датчики дотику, а після легкого ушкодження стають навіть чутливішими~\cite{LaMotte1982hyper}; послаблення відповіді після сильного нагрівання виміряно лише в одного їхнього типу~\cite{Treede1995heat}. Датчик світла на сталому тлі замовкає (рис.~\ref{fig:adaptation}), а сигналізація, яка замовкла через хвилину, марна. По-друге, вона має високий поріг: датчики болю спрацьовують лише тоді, коли дія близька до небезпечної, наприклад на нагрівання --- у різних датчиків пороги від $41^\circ$ до $46$--$48^\circ$, залежно від типу~\cite{Treede1995heat, Weidner1999cfib}. По-третє --- і це головне для цього розділу --- гучність тривоги вирішує не лише датчик, а й сама машина. Та сама рана болить по-різному в різних обставинах. Подивімося, з яких уже знайомих деталей це зібрано.

Датчики болю --- \nt{GLUT}-нейрони, як і решта чутливих нейронів розділу~\ref{sec:sensors}. Частина з них разом із глутаматом виділяє речовину~P з додатка~\ref{sec:othermediators}, яка повільно підсилює передачу.

\section{Два проводи: швидкий і повільний біль}

Ударте палець --- і ви відчуєте біль двічі: спершу короткий і різкий, а через секунду тупий і довгий. Це два різні проводи. Швидкий провід ізольований і проводить імпульси зі швидкістю $v$ від $5$ до $30$ метрів за секунду; повільний тонкий і голий, $0.5$--$2$ метри за секунду. Сигнал від стопи проходить близько метра, і час приходу
\begin{empheq}[box=\keybox]{equation}
t \;=\; \frac{L}{v}
\label{eq:paintime}
\end{empheq}
для швидкого проводу --- десятки мілісекунд, для повільного --- близько секунди. Два проводи несуть два повідомлення. Швидкий каже \emph{де} і \emph{коли}: його імпульси приходять раніше й точніше, як у коді часом розділу~\ref{sec:code}. Повільний каже, \emph{наскільки погано}, і його довгий тупий біль тримає машину в режимі тривоги, доки ушкодження не минуло.

\section{Ворота в спинному мозку}

Перше, куди надходить сигнал болю, --- спинний мозок, і там він проходить через \emph{ворота}~\cite{MelzackWall1965}; конкретні нейрони цих воріт знайдено порівняно недавно~\cite{Duan2014}. На вихідний \nt{GLUT}-нейрон, що передає біль далі в мозок, сходяться провід болю й провід дотику. А поруч сидить гальмівний нейрон-посередник, \nt{GABA} або \nt{GLYC}, якого дотик збуджує (рис.~\ref{fig:gate}). Дотик зачиняє ворота. У вихідній теорії воріт~\cite{MelzackWall1965} біль, навпаки, вимикає цього посередника, і сильний біль ворота відчиняє; це припущення теорії, на знайдених нейронах воріт його прямо не показано.

\begin{figure}[t]
\centering
\begin{tikzpicture}[>=Stealth,
  nrn/.style={circle,draw,very thick,minimum size=1.0cm,inner sep=0pt,font=\scriptsize},
  inh/.style={-{Bar[width=7pt]},thick,red!70!black},
  bc/.style={->,thick,densely dashed,orange!85!black}]
\node[nrn,fill=blue!6] (Z) at (6,0) {\nt{GLUT}};
\node[nrn,fill=red!10] (Q) at (3.6,-1.6) {\nt{GABA}};
\draw[->,very thick,red!70!black] (0,0.6) node[left,font=\small,black,align=right] {біль $\nu$} -- (5.45,0.15);
\draw[->,very thick,blue!60!black] (0,-2.6) node[left,font=\small,black,align=right] {дотик $\mu$} -- (2.9,-2.0);
\draw[->,thick,blue!60!black] (1.8,-2.35) -- (5.5,-0.35) node[pos=0.8,below right=-2pt,font=\scriptsize] {$w_{z\mu}$};
\draw[inh] (1.6,0.5) -- (3.35,-1.1) node[pos=0.5,left=1pt,font=\scriptsize] {$w_{q\nu}$};
\draw[inh] (Q) -- (Z) node[pos=0.5,above left=-1pt,font=\scriptsize] {$\beta$};
\draw[->,very thick] (Z) -- (8.2,0) node[right,font=\small,align=left] {у мозок: $z$};
\draw[bc] (6,2.1) node[above,font=\scriptsize,black,align=center] {згори: \nt{SERO}, \nt{NORA}, опіоїди} -- (Z.north) node[pos=0.45,right,font=\scriptsize,black] {підсилення $g$};
\end{tikzpicture}
\caption{Ворота болю в спинному мозку за моделлю~\eqref{eq:gate}. Вихідний \nt{GLUT}-нейрон отримує біль $\nu$ і слабке збудження від дотику $\mu$. Гальмівний посередник (\nt{GABA} або \nt{GLYC}) збуджується дотиком і, за припущенням теорії воріт, вимикається болем. Згори, широкомовними каналами, надходить підсилення $g$.}
\label{fig:gate}
\end{figure}

Запишімо це для частот, як у розділі~\ref{sec:gaba}. Частота посередника $q$ і вихідна частота $z$ дорівнюють
\begin{empheq}[box=\keybox]{equation}
q \;=\; \bigl[w_{q\mu}\,\mu + q_0 - w_{q\nu}\,\nu\bigr]_+ ,
\qquad
z \;=\; \bigl[\,g\,(\nu + w_{z\mu}\,\mu) - \beta\,q\,\bigr]_+ ,
\label{eq:gate}
\end{empheq}
де $\nu$ --- вхід болю, $\mu$ --- вхід дотику, $w$ --- ваги зв'язків (перший індекс --- кому, другий --- від кого), $\beta$ --- сила гальмування, $q_0$ --- власна фонова активність посередника, $g$ --- підсилення, яке задають згори (про нього нижче). Це заборона~\eqref{eq:veto} з розділу~\ref{sec:gaba}: «біль, але не за дотику», --- з однією поправкою, взятою з теорії воріт як припущення: сильний біль сам вимикає забороняльний нейрон.

Ми порахували модель при $w_{q\mu}=1$, $q_0=0.2$, $w_{q\nu}=0.5$, $w_{z\mu}=0.3$, $\beta=1.5$ (рис.~\ref{fig:gatecurves}). Без дотику біль починає проходити з $\nu\approx0.18$. З дотиком поріг піднімається до $0.86$, і при $\nu=1$ вихід падає вчетверо, з $1$ до $0.25$. А за сильного болю, $\nu=2$, дотик уже нічого не змінює: ворота розчахнуті. Так і в житті: потерти забите місце допомагає, а за серйозної травми --- ні. Сам дотик, без болю, через ворота не проходить: тривога від дотику не здіймається.

\begin{figure}[t]
\centering
\begin{tikzpicture}[>=Stealth,x=5.2cm,y=1.35cm]
\draw[->,gray!70!black] (0,0) -- (2.12,0) node[right,font=\small,black] {біль $\nu$};
\draw[->,gray!70!black] (0,0) -- (0,4.3) node[left,font=\small,black] {$z$};
\foreach \t in {0,0.5,1,1.5,2}{\draw[gray!60!black] (\t,-0.05) -- (\t,0.05); \node[below,font=\scriptsize] at (\t,0) {\t};}
\foreach \f in {1,2,3,4}{\draw[gray!60!black] (-0.01,\f) -- (0.01,\f); \node[left,font=\scriptsize] at (0,\f) {\f};}
\draw[very thick,red!70!black] plot file {figures/pain_gate_notouch.dat};
\draw[very thick,blue!60!black] plot file {figures/pain_gate_touch.dat};
\draw[very thick,orange!85!black,densely dashed] plot file {figures/pain_gate_sens.dat};
\draw[very thick,red!70!black] (0.08,3.9) -- (0.2,3.9) node[right,font=\scriptsize,black] {без дотику};
\draw[very thick,blue!60!black] (0.08,3.45) -- (0.2,3.45) node[right,font=\scriptsize,black] {з дотиком};
\draw[very thick,orange!85!black,densely dashed] (0.08,3.0) -- (0.2,3.0) node[right,font=\scriptsize,black] {після сенсибілізації, $g=2$};
\end{tikzpicture}
\caption{Вихід воріт~\eqref{eq:gate} залежно від сили болю. Дотик піднімає поріг і приглушує слабкий і середній біль, але не сильний. Сенсибілізація (штрихова лінія) вдвічі збільшує підсилення: той самий біль передається вдвічі гучніше, а поріг опускається.}
\label{fig:gatecurves}
\end{figure}

\section{Ручка гучності згори}

Ворота керуються не лише знизу. Зі стовбура мозку до них ідуть низхідні шляхи, які змінюють підсилення $g$ у~\eqref{eq:gate}: \nt{SERO}, \nt{NORA} та опіоїдні пептиди з додатка~\ref{sec:othermediators}~\cite{Fields2004}. Це ті самі широкомовні канали, що й у розділах~\ref{sec:serocomputer}--\ref{sec:acet}, тільки тут їхня ручка --- гучність тривоги. Вони вміють крутити її в обидва боки: ті самі низхідні кола можуть і приглушувати біль, і посилювати його.

Навіщо машині зменшувати тривогу? Тому що тривога --- це переривання, а переривання не завжди доречне. Тварина, що тікає від хижака, не повинна зупинятися через поранену лапу: спершу тікати, потім зализувати рану. Очікування полегшення теж приглушує біль, і частина цього ефекту зникає, якщо заблокувати опіоїдні рецептори~\cite{Levine1978}: очікування, обчислене в мозку, спускається опіоїдним каналом до воріт. Саму цю картину виміряно; як саме мозок вирішує, якою має бути гучність, --- гіпотеза.

\section{Сенсибілізація: тривога, що навчається}

Після ушкодження вихідні нейрони воріт стають чутливішими: той самий вхід викликає більший вихід, а поріг опускається~\cite{Woolf1983}. У моделі~\eqref{eq:gate} це зростання підсилення $g$, і найпростіший його опис --- повільне RC-коло, яке накачує сам біль:
\begin{empheq}[box=\keybox]{equation}
\tau_s\,\frac{\dd g}{\dd t} \;=\; -(g-1) + k_s\,z ,
\label{eq:sensitization}
\end{empheq}
де $\tau_s$ --- час, за який чутливість повертається до норми; він більший за тривалість самого болю, бо сенсибілізація тримається й після того, як ушкоджувальний стимул припинився~\cite{Woolf1983}, а $k_s$ --- сила накачування. Доки тканину ушкоджено, вихід воріт тримає $g$ вище від одиниці, і все, що відбувається поруч із раною, передається гучніше (штрихова лінія на рис.~\ref{fig:gatecurves}: при $g=2$ поріг опускається до $0.11$, а вихід подвоюється). Це розумне налаштування сигналізації: доки ушкодження не загоїлося, краще перестрахуватися.

Для інженера це додатний зворотний зв'язок із повільним витоком: біль підвищує підсилення, підсилення --- біль. Доки петля слабка, $g$ повертається до норми, коли рана загоїлася. Якщо ж петля сильна, сигналізація може застрягнути в увімкненому стані й після того, як ушкодження вже немає, --- як група із сильним зворотним зв'язком на рис.~\ref{fig:bistable}. Модель~\eqref{eq:sensitization} --- спрощення; те, що центральна сенсибілізація існує, виміряно.

\section{Біль як учитель і як переривання}

Біль у машині має дві роботи, крім сигналізації.

\emph{Учитель.} Біль --- від'ємна винагорода: у формулу помилки~\eqref{eq:ctd} він входить як $r<0$. Усе, що говорилося в розділі~\ref{sec:dofa}, працює й тут: провісник болю починає викликати помилку заздалегідь, і мережа вчиться уникати того, що до болю веде. Досліди на людях показують, що активність мозку під час навчання на болю йде саме за помилкою часових різниць~\cite{Seymour2004}, а частина \nt{DOPA}-нейронів відповідає й на неприємні події (розділ~\ref{sec:dofaalt}).

\emph{Переривання.} Біль відриває увагу від поточної задачі --- це його призначення, а не побічний ефект~\cite{EcclestonCrombez1999}. У розділі~\ref{sec:nora} ту саму роль «переривання» обговорювали для уривчастої пачки \nt{NORA}. А найшвидша відповідь навіть не доходить до мозку: відсмикування руки від гарячого замикається просто в спинному мозку тією самою парою м'язів і тією самою забороною антагоніста, що на рис.~\ref{fig:antagonists}.

\begin{keyblock}
\begin{proposition}[Сигнал тривоги]
\label{prop:pain}
Біль --- не вимірювання, а сигнал тривоги з високим пріоритетом. Він звикає набагато слабше за вимірювальні канали, іде двома проводами (швидким --- «де», повільним --- «наскільки погано»), проходить через ворота, які зачиняє дотик (а сильний біль, за припущенням теорії воріт, відчиняє), і його гучність задають згори широкомовні канали. Ці пристрої виміряно, крім відчинення воріт болем; прості моделі~\eqref{eq:gate} і~\eqref{eq:sensitization} --- спрощення.
\end{proposition}
\end{keyblock}

\section{Підсумок}

\begin{table}[t]
\centering
\footnotesize
\setlength{\tabcolsep}{4pt}
\renewcommand{\arraystretch}{1.15}
\begin{tabular}{>{\raggedright}p{3.0cm}>{\raggedright}p{4.9cm}>{\raggedright\arraybackslash}p{4.9cm}}
\toprule
Що робить біль & Як & Що відомо \\
\midrule
здіймає тривогу & високий поріг, звикання слабше, ніж у дотику & поріг виміряно; порівняння звикання --- оцінка \\
повідомляє «де» і «наскільки погано» & два проводи з різною швидкістю~\eqref{eq:paintime} & виміряно \\
проходить через ворота & заборона через \nt{GABA}/\nt{GLYC}~\eqref{eq:gate} & ворота виміряно; відчинення болем --- припущення теорії; модель --- спрощення \\
змінює гучність & низхідні \nt{SERO}, \nt{NORA}, опіоїди & виміряно; правило вибору гучності --- гіпотеза \\
навчається після ушкодження & зростання підсилення~\eqref{eq:sensitization} & сенсибілізацію виміряно; модель --- спрощення \\
навчає уникати & від'ємна винагорода в~\eqref{eq:ctd} & схожість із помилкою часових різниць виміряно \\
перериває роботу & захоплення уваги; рефлекс відсмикування & виміряно \\
\bottomrule
\end{tabular}
\caption{Біль як сигнал тривоги.}
\label{tab:pain}
\end{table}

Біль зібрано з деталей, які ми вже бачили (таблиця~\ref{tab:pain}): \nt{GLUT}-датчики, два проводи, заборона через гальмівного посередника, широкомовна ручка гучності, повільне RC-коло, помилка передбачення й переривання. Нове в ньому одне: це єдиний вхід машини, гучність якого машина задає сама, --- сигналізація, яку власник може і приглушити, і переналаштувати.

\section{Задачі}

\begin{exercise}[\lvlA]
\label{ex:14:1}
\emph{Два проводи.} Сигнал іде від стопи, $L=1$~м. (а)~Залп іде пучком проводів одного типу, швидкості яких заповнюють діапазон~\eqref{eq:paintime}. На скільки він розтягується до приходу в спинний мозок для швидкого й повільного типу? (б)~Хвилі болю проводами $15$ і $1$~м/с розрізненні за різниці від $0.1$~с. З якої відстані вони зливаються?
\end{exercise}

\begin{exercise}[\lvlB]
\label{ex:14:2}
\emph{Коли дотик болить.} Ворота~\eqref{eq:gate} з параметрами тексту, підсилення $g$ зростає під час сенсибілізації. До якого $g$ дотик без болю не проходить через ворота за жодної сили $\mu$? За якого $g$ починає проходити дотик $\mu=1$?
\end{exercise}

\begin{exercise}[\lvlB]
\label{ex:14:3}
\emph{Стійкість сенсибілізації.} Входи сталі, вихід воріт лінійний за $g$: $z=gA-C$, $A=\nu+w_{z\mu}\mu$, $C=\beta q\ge0$. (а)~Знайдіть рівновагу $g^*$ моделі~\eqref{eq:sensitization} і умову її стійкості. (б)~При $k_s=0.5$ знайдіть силу болю, вище від якої модель іде врознос, без дотику і з дотиком $\mu=1$.
\end{exercise}

\begin{exercise}[\lvlA]
\label{ex:14:4}
\emph{Провісник болю.} Сигнал у момент $0$ передбачає біль у момент $T=2$~с, у~\eqref{eq:ctd} $r(t)=-P\,\delta(t-T)$, $\tau_\gamma=2$~с. (а)~Знайдіть площу сплеску $\delta$ у момент сигналу. (б)~Дія в момент $t_e=1$~с скасовує біль. Який сплеск $\delta$ у цей момент і чого він навчить мережу?
\end{exercise}

\begin{exercise}[\lvlB]
\label{ex:14:5}
\emph{Сигналізація, що защіпається.} Доповніть ворота з \texttt{models/\allowbreak pain\_gate.py} динамікою~\eqref{eq:sensitization} і насиченням $z\le4$. Дотик $\mu=1$ є завжди, ушкодження $\nu=2$ триває час $T$, далі $\nu=0$. Моделюванням знайдіть найменше $T$ (в одиницях $\tau_s$), після якого сигналізація лишається ввімкненою, при $k_s=4$, $4.6$, $5$, $6$, $8$.
\end{exercise}

\begin{exercise}[\lvlB]
\label{ex:14:6}
\emph{Проєктування воріт.} При $\beta=1.5$, $w_{z\mu}=0.3$, $g=1$ доберіть у~\eqref{eq:gate} $q_0$, $w_{q\nu}$, $w_{q\mu}$ так, щоб поріг болю був $0.2$ без дотику і $1.0$ з дотиком $\mu=1$, а дотик переставав приглушувати біль при $\nu_\times=2.5$. До якого підсилення $g$ дотик без болю не проходить через ворота?
\end{exercise}

\begin{exercise}[\lvlC]
\label{ex:14:7}
\emph{Ручка гучності.} Робота приносить цінність $V$ за одиницю часу, робота з ушкодженням $\nu$ коштує $c\nu$, тож переривати вигідно при $\nu>V/c$; біль перериває при $z>\theta=0.5$. (а)~Яке підсилення $g^*$ воріт~\eqref{eq:gate} без дотику дає цей поріг при $V/c=0.25$, $0.5$, $2$? (б)~З яких сигналів книги машина могла б оцінювати $V$ і $c$?
\end{exercise}
